\chapter{Medical Implants}

\medterm{Artificial Urinary Sphincter (AUS)}-- A fully implantable device for treating severe stress urinary incontinence, most commonly in men after radical prostatectomy. Components: an inflatable periurethral cuff, a pressure-regulating balloon (placed in the space of Retzius), and a control pump (placed in the scrotum/labium). The patient activates the pump to deflate the cuff, allowing voiding; the cuff automatically refills over 3-5 minutes. The AMS 800 (Boston Scientific) is the most widely used device. Complications: erosion (cuff into urethra, 5-10\% at 10 years), infection, mechanical failure, and persistent incontinence. Revision rates: 25-40\% at 10 years.

\medterm{Artificial Ventricle (Total Artificial Heart)}-- A pulsatile, pneumatically or electrically driven biventricular replacement that removes the native ventricles. The SynCardia temporary Total Artificial Heart (TAH-t) is FDA-approved as a bridge to cardiac transplantation. The Carmat TAH is a bioprosthetic option available in Europe. Indications: severe biventricular failure, massive myocardial infarction, and end-stage heart disease when a donor heart is unavailable. Implantation involves full median sternotomy, cardiopulmonary bypass, and excision of the native ventricles. Pneumatic drivers (Big Blue, Companion 2) power the system externally. Complications: stroke, bleeding, infection (mediastinitis, driveline site), hemolysis, and device malfunction. Longest survival: >4 years as bridge to transplant.

\medterm{BAC (Bone Anchored Cochlear) -- See Bone-Anchored Hearing Aid (BAHA)}-- An auditory implant that uses a titanium screw implanted in the temporal bone behind the ear with a sound processor attached. The implant bypasses the outer and middle ear by conducting sound through bone directly to the cochlea. Widely known as BAHA (Bone-Anchored Hearing Aid). The modern term is BCI (Bone Conduction Implant).

\medterm{BAHA (Bone-Anchored Hearing Aid)}-- A surgically implantable bone conduction hearing system. Components: a titanium abutment (fixture) placed in the mastoid bone, and an external sound processor that clips onto the abutment. The fixture osseointegrates with bone (takes 3-12 weeks). Indications: conductive hearing loss (chronic otitis media, aural atresia, otosclerosis), mixed hearing loss, and single-sided deafness (SSD). Contraindications: poor bone quality, uncontrolled diabetes, prior radiation to the temporal bone, and age <5 years. The Cochlear Baha system and Oticon Medical Ponto are the major brands. Complications: skin infection (peri-abutment dermatitis, most common), fixture failure (loss of osseointegration), and soft tissue overgrowth. Implant survival: >95\% at 5 years.

\medterm{Biodegradable Stent (Bioresorbable Stent)}-- A coronary stent made of a bioabsorbable material (poly-L-lactic acid, magnesium alloy) that gradually dissolves over 2-3 years, theoretically restoring vascular vasomotion and avoiding the long-term risks of permanent metallic stents (e.g., late stent thrombosis, neoatherosclerosis). The Absorb BVS (Abbott) was the first FDA-approved bioresorbable stent but was withdrawn due to higher target lesion failure and scaffold thrombosis rates compared to drug-eluting stents. Second-generation scaffolds (e.g., Magmaris magnesium stent, Fantom sirolimus-eluting scaffold) aim to improve outcomes. Complications: scaffold thrombosis (higher than DES, especially in small vessels), malapposition, and fracture. Current use: selective, in clinical trials or specific anatomical settings.

\medterm{Bone Graft (Implantable)}-- Natural or synthetic bone material implanted to promote bone healing, fill defects, or achieve spinal fusion. Types: autograft (from the patient — iliac crest, fibula, rib, distal radius — the gold standard for osteogenesis, osteoconduction, and osteoinduction), allograft (cadaveric — cortical strut, cancellous chips, demineralized bone matrix DBM), synthetic (calcium phosphate hydroxyapatite, tricalcium phosphate, bioactive glass, calcium sulfate), and growth factor-enhanced (BMP-2 on an absorbable collagen sponge, e.g., Infuse Bone Graft). Formulations: putty, paste, granules, sheets, and structural blocks. Autograft limitations: donor site morbidity (pain, infection, hematoma, fracture) and limited quantity. Allograft risks: infection (very low with modern processing), and fracture. DBM: osteoinductive but inconsistent quality. BMP-2: potent but associated with ectopic bone formation, radiculitis, and cancer risk.

\medterm{Bone Plate and Screws--(Implantable Fracture Fixation)}-- Metal plates and screws used to stabilize bone fractures or osteotomies. Materials: stainless steel (strong, cheap, corrodes), titanium (biocompatible, MRI-safe, elastic modulus closer to bone), and bioabsorbable (PLLA, PGA — for children, no removal needed). Locking vs. non-locking screws: locking screws thread into the plate, creating a fixed-angle construct (like an external fixator internally), ideal for osteoporotic bone. Compression plating: dynamic compression plate (DCP) and limited-contact DCP (LC-DCP) for anatomic reduction and axial compression. Neutralization plate: protects lag screws from rotational forces. Buttress plate: supports a metaphyseal fragment (e.g., tibial plateau). Bridge plate: spans a comminuted fracture zone without direct reduction. Anatomical plates: pre-contoured for specific bones (distal radius, proximal humerus, acetabulum, calcaneus). The AO/ASIF principles guide fixation: anatomic reduction (articular), stable fixation, preservation of blood supply, and early mobilization. Complications: nonunion/malunion, infection (hardware removal if chronic), screw breakage, plate loosening, stress shielding (bone resorption under rigid plates), and refracture after removal.

\medterm{Breast Implant (Prosthesis)}-- A silicone- or saline-filled device placed in the subglandular or submuscular plane for breast augmentation (cosmetic) or reconstruction (post-mastectomy). Silicone gel implants: cohesive gel ("gummy bear") that retains shape even if the shell ruptures. Saline implants: filled intraoperatively, less natural feel. Surface types: smooth (less friction, higher capsular contracture) and textured (lower capsular contracture, but associated with BIA-ALCL — breast implant-associated anaplastic large cell lymphoma). Shape: round (most common) and anatomical (teardrop, requires textured surface to prevent rotation). FDA-approved manufacturers: Allergan (Natrelle), Mentor (MemoryGel, MemoryShape), Sientra, and Motiva (ergonomix). Implant placement: subglandular (on top of pectoralis, more animation deformity), subpectoral (dual plane, partial under muscle), and prepectoral (on top of muscle, for reconstruction). Complications: capsular contracture (most common, Baker grade I-IV), rupture (silent in silicone gel, MRI recommended at 5-6 years post-op then q2-3 years), implant malposition, rippling, infection, hematoma, and BIA-ALCL (textured implants, 1:3,000 to 1:30,000). Implant life: 10-20 years (not lifetime devices).

\medterm{Cochlear Implant}-- An electronic medical device that bypasses damaged hair cells in the cochlea and directly stimulates the auditory nerve. Components: internal (receiver-stimulator placed in a bony well behind the ear, electrode array inserted into the scala tympani of the cochlea) and external (microphone, speech processor worn behind the ear or body-worn, and transmitter coil held in place by a magnet). Indications: severe-to-profound bilateral sensorineural hearing loss with insufficient benefit from hearing aids (aided word recognition <50-60\% in the best-aided condition). Adult candidacy: post-lingual onset. Pediatric candidacy: FDA-approved down to 9 months of age; earlier implantation improves language outcomes. Major brands: Cochlear (Nucleus series), Advanced Bionics (HiRes Ultra, Naída CI), and MED-EL (SYNCHRONY, RONDO). Speech processors replaceable/upgradable externally. Surgery: mastoidectomy, facial recess approach, cochleostomy or round window insertion, NRT/telemetry to confirm placement. Complications: facial nerve injury (rare, <1\%), meningitis (pneumococcal vaccination required pre-op), flap infection, electrode migration, device failure (1-3\%), and vestibular symptoms (vertigo). Outcomes: most achieve open-set speech understanding; factors are age at implantation, duration of deafness, and adherence to aural rehabilitation.

\medterm{Continuous Glucose Monitor (CGM Implant)}-- A subcutaneous sensor that measures interstitial glucose levels every 5-15 minutes. The Dexcom G7 and G6 (FDA-approved for non-adjunctive use, i.e., can dose insulin without fingersticks), Abbott FreeStyle Libre (flash glucose monitoring, reader on demand with 8-hour continuous data), and Medtronic Guardian (integrated with insulin pumps). Inserted subcutaneously in the abdomen (or arm for Libre) and worn for 7-14 days. Fully implantable CGM (Eversense, Senseonics) with a subcutaneous fluorescence-based sensor lasting 90-180 days; requires surgical insertion in an office procedure. MARD (mean absolute relative difference): <10\%. Alarms: hypo/hyperglycemia, rate of change, and predictive alerts. Benefits: reduced HbA1c, less hypoglycemia, improved time-in-range (TIR, 70-180 mg/dL), and integration with automated insulin delivery (AID) systems. Limit: lag behind blood glucose (5-10 minutes interstitial delay).

\medterm{Coronary Stent}-- A metallic mesh tube deployed in a coronary artery to treat atherosclerotic narrowing. Bare metal stent (BMS): first generation, high restenosis rate (20-40\%) due to neointimal hyperplasia. Drug-eluting stent (DES): releases an antiproliferative drug (sirolimus, paclitaxel, everolimus, zotarolimus) from a polymer coating to inhibit restenosis. Second-generation DES (everolimus: Xience, Promus; zotarolimus: Resolute): thinner struts, biocompatible polymers, 5-10\% restenosis, and lower late stent thrombosis (LST) than first-generation (Cypher, Taxus). Third-generation: bioabsorbable polymer (Orsiro, Synergy) or polymer-free (BioFreedom) to reduce inflammation. Ultra-thin strut DES (<65 µm): improved deliverability. Cobalt-chromium and platinum-chromium alloys provide radial strength with thinner struts. Current standard: DES is used in >95\% of PCI. Dual antiplatelet therapy (DAPT: aspirin + P2Y12 inhibitor) is required after deployment to prevent in-stent thrombosis, typically for 3-6 months (new-generation DES) to 12 months (high-risk ACS), then aspirin lifelong. In-stent restenosis (ISR): treated with drug-coated balloon (SeQuent Please, Lutonix) or another DES. Very late stent thrombosis (>1 year): rare with modern DES (0.2-0.5\%/year).

\medterm{Deep Brain Stimulator (DBS)}-- A neurostimulator that delivers electrical pulses to targeted deep brain structures via implanted electrodes. Components: intracranial leads (quadripolar electrodes) placed stereotactically, extension cables tunneled subcutaneously, and an implantable pulse generator (IPG, "battery") placed in the infraclavicular or abdominal subcutaneous pocket. Targets: subthalamic nucleus (STN — for Parkinson's disease, most common), globus pallidus internus (GPi — for Parkinson's and dystonia), ventral intermediate nucleus (VIM — for essential tremor), and anterior nucleus of the thalamus (ANT — for epilepsy). The Medtronic Percept PC (with sensing capability, BrainSense), Abbott Infinity (directional leads), and Boston Scientific Vercise (multiple independent current control) are current IPG platforms. DBS is MRI-conditional (specific conditions). Stimulation parameters: amplitude (voltage), pulse width (60-90 µs), frequency (130-180 Hz for Parkinson's), and active electrode configuration. Indications: Parkinson's disease (motor fluctuations, dyskinesias, tremor refractory to medication), essential tremor, dystonia (primary generalized), obsessive-compulsive disorder (OCD, under HDE), and epilepsy. Surgery: frame-based (CRW, Leksell) or frameless (NextFrame, StealthStation) stereotaxy, microelectrode recording (MER) for physiological mapping, intraoperative test stimulation for symptom improvement and side effect assessment. Complications: hemorrhage (1-2\%, intraparenchymal or along the lead tract), infection (2-5\%, often IPG pocket), lead fracture/migration, hardware erosion, stimulation-induced side effects (paresthesias, dysarthria, mood changes), and IPG battery depletion (replaced every 2-5 years, rechargeable IPGs are emerging).

\medterm{Dental Implant (Endosseous)}-- A titanium or titanium-alloy screw-shaped post that osseointegrates with the alveolar bone to support a crown, bridge, or denture. Materials: commercially pure titanium (Grade 4, rough surface for osseointegration) and zirconia (esthetic, metal-free, but less evidence for posterior teeth). Two-stage: implant placed, submerged under gingiva for 3-6 months of osseointegration, then uncovered and loaded. Single-stage: transgingival healing abutment placed at the time of implant insertion, avoiding a second surgery. Immediate loading: crown placed within 48 hours (for anterior teeth, ideal bone quality). Delayed loading: standard after 3-6 months. The Branemark protocol (1965) established the principles: atraumatic extraction, sterile surgical technique, and uninterrupted healing. Brands: Straumann (SLA, Roxolid), Nobel Biocare (Replace, Brånemark), Zimmer (Tapered Screw-Vent), Dentsply (Astra Tech, OsseoSpeed), and BioHorizons. Surface modifications: SLA (sandblasted + acid-etched), TiUnite (oxidized), and SLActive (hydrophilic, faster osseointegration). Contraindications: uncontrolled diabetes, bisphosphonate therapy (MRONJ risk), heavy smoking, and history of head/neck radiation. Complications: peri-implantitis (inflammation and bone loss, 20-30\% of patients), implant failure (early: failed osseointegration, 1-5\%; late: due to overloading or infection), nerve injury (inferior alveolar nerve, mental nerve), sinus perforation (maxillary sinus), and esthetic problems (gum recession visible metal margin). 10-year survival: >95\% in well-selected patients.

\medterm{Glaucoma Drainage Device (Tube Shunt)}-- An implant for reducing intraocular pressure (IOP) in glaucoma by draining aqueous humor from the anterior chamber to an equatorial bleb. Components: a silicone tube (placed in the anterior chamber through a scleral track) connected to an end plate (positioned under the rectus muscles on the scleral surface). The Ahmed valve (New World Medical) has a flow restrictor (venturi-based or silicone membrane) that limits hypotony immediately post-op. The Baerveldt implant (Johnson \& Johnson) has no valve; flow is restricted by a ligation suture or dissolvable Vicryl tie for a 4-6 week delay. The Molteno implant (single- or double-plate) is the original. Indications: neovascular glaucoma, uveitic glaucoma, post-penetrating keratoplasty glaucoma, prior failed trabeculectomy, congenital glaucoma, and high-risk glaucoma. Contraindications: corneal endothelial disease (tube can cause corneal decompensation) and active infection. Surgery: superotemporal quadrant, peritomy, quadrant dissection, tube trimmed bevel-up, scleral patch graft (donor sclera, pericardium, or cornea) over the tube entry site, and conjunctival closure. Complications: hypotony (especially Baerveldt without ligation), tube erosion (conjunctival breakdown), corneal edema, tube obstruction (iris, fibrin, vitreous), choroidal effusion, endophthalmitis, and diplopia from the plate interfering with extraocular muscles. Tube shunt survival: >80\% at 2 years, decreasing to about 50\% at 5 years.

\medterm{Implantable Cardioverter-Defibrillator (ICD)}-- An electronic device that monitors heart rhythm and delivers therapy (antitachycardia pacing ATP, cardioversion, or defibrillation) for ventricular arrhythmias. Components: pulse generator (placed in the left infraclavicular subcutaneous pocket) and one or more leads (transvenous: single-coil or dual-coil, placed in the RV apex or septum; or entirely subcutaneous S-ICD). Indications (primary prevention): ischemic cardiomyopathy (EF \(\le\)35\% >40 days post-MI), non-ischemic cardiomyopathy (EF \(\le\)35\%, on GDMT for \(\ge\)3 months), and certain channelopathies (long QT, Brugada, HCM, ARVC). Secondary prevention: survivors of SCD due to VF/VT. Transvenous ICD: pectoral implant, 1-3 leads (RA, RV, LV for CRT-D). Subcutaneous ICD (S-ICD, Boston Scientific Emblem): no intravascular leads, left parasternal electrode and left lateral pocket, for patients without need for pacing/bradycardia pacing. Leadless ICD: entirely within the RV (currently only a pacemaker is available commercially leadless). Shock: up to 41 J delivered biphasic waveform. Defibrillation threshold (DFT): <10 J difference from maximum output is acceptable. Programming: optimizing VT detection zones (monitor, ATP, shock), S-ICD with SMART Pass algorithm reduces inappropriate shocks. Complications: infection (pocket, endocarditis, highest with system revision), lead failure (fracture, insulation break), inappropriate shocks (SVT, AF, lead noise, T-wave oversensing), pneumothorax/hemothorax (transvenous), and device erosion. Device generators last 5-10 years (transvenous), longer with S-ICD (no pacing drain). Remote monitoring is standard (CareLink, Merlin, Latitude) and reduces clinic visits and improves survival.

\medterm{Implantable Loop Recorder (ILR)}-- A subcutaneous cardiac monitoring device that continuously records a single-lead ECG for up to 3 years. The size of a thumb drive, inserted in the left parasternal area (4th intercostal space) in a 1-cm pocket under local anesthesia. The Medtronic LINQ II (micra-like, injectable), Abbott Confirm Rx (Bluetooth-enabled, smartphone-based transmission), and Biotronik BioMonitor 3 are current models. Autotriggered recordings: asystole (>3 sec), bradycardia, tachycardia, and AF detection (P-wave analysis, AF burden\%). Patient-activated recordings: for symptoms (palpitations, syncope, presyncope, chest pain). Indications: syncope of unclear etiology (after negative workup — most common indication), cryptogenic stroke (to detect occult AF, 20-30\% detection rate), and AF management (AF burden monitoring after ablation or cardioversion). Ill at TIA/CVA workup: monitor 3 years, detects AF in up to 20-30\% of patients. Battery life: 2-3 years. The device is explained at the end of monitoring or battery depletion. Complications: infection (1-3\%), pocket hematoma, device migration, and discomfort.

\medterm{Insulin Pump (CSII -- Continuous Subcutaneous Insulin Infusion)}-- A battery-powered device that delivers rapid-acting insulin (analog) subcutaneously through a cannula, providing a continuous basal rate and patient-activated boluses. Types: tethered pump (Medtronic MiniMed 780G, Tandem t:slim X2 — with touchscreen, tubed, Control-IQ/Advanced Hybrid Closed Loop), patch pump (Omnipod DASH/Omnipod 5 — tubeless, fully disposable, worn for 3 days, controlled via a PDM or smartphone), and the Ypsomed mylife YpsoPump (tubed, color touchscreen). Hybrid closed-loop (HCL, automated insulin delivery AID): integrates CGM data to automatically adjust basal insulin; the Medtronic 780G (SmartGuard, targets BG 100-120 mg/dL) and Tandem Control-IQ (target 112.5-160 mg/dL) are FDA-approved. The Omnipod 5 closed-loop is the first tubeless HCL. Indications: type 1 diabetes (all patients considered), poorly controlled type 2 diabetes, frequent/severe hypoglycemia, dawn phenomenon, and gastroparesis. Cannula: inserted subcutaneously (abdomen, arm, hip, thigh) with an automatic inserter, changed every 2-3 days. Insulin: only rapid-acting (Humalog, Novolog, Fiasp, Lyumjev). Reservoirs: 180-300 units. Contraindications: unwillingness/unable to manage the system, untreated eating disorder, and severe visual impairment without caregiver support. Complications: infusion set failure (kinked cannula, dislodgment, occlusion, bleeding) leading to DKA within 4-6 hours (no long-acting insulin on board), lipohypertrophy, infection at the insertion site, and pump malfunction.

\medterm{Intraocular Lens (IOL)}-- An artificial lens implanted in the eye after cataract extraction (phacoemulsification). One-piece acrylic (monofocal, aspheric, with blue-blocking chromophore) is the standard of care placed in the capsular bag. Monofocal IOL: single focal point (set for distance), requires reading glasses. Multifocal IOL: diffractive rings on the surface divide light into multiple focal points (distance, near, intermediate), reduces spectacle dependence (e.g., Alcon Restor, Johnson \& Johnson Tecnis Symfony). Depth-of-focus (EDOF) IOL: extends the focal range with less dysphotopsia than multifocal (e.g., Tecnis Symfony, Alcon Vivity). Toric IOL: corrects pre-existing corneal astigmatism (\(\ge\)0.75 D); requires careful alignment (axis marking). Accommodating IOL: changes shape with ciliary muscle contraction (e.g., Crystalens, no longer commonly used). Light-adjustable lens (LAL, RxSight): can be adjusted postoperatively with UV light to fine-tune refractive outcome. Blue-blocking IOLs: filter UV and blue light (theoretical protection against macular degeneration, though controversial). Biometry: optical (IOLMaster, Lenstar) measures axial length, K-readings, and ACD; formulas (SRK/T, Holladay 2, Barrett Universal II, Hill-RBF) calculate IOL power. Target refraction: 0.0 D (plano) for distance, -0.25 D to -0.50 D for most patients. Contraindications: zonular weakness (pseudoexfoliation, trauma), capsular tear, and uveitis (active). Complications: posterior capsule opacification (PCO, 20-40\% at 5 years, treated with YAG capsulotomy), IOL dislocation/decentration, cystoid macular edema (Irvine-Gass syndrome), dysphotopsia (positive: arc-like light streaks; negative: temporal shadow), and refractive error. Phakic IOL (ICL, implantable collamer lens): placed in the sulcus (EVO ICL by STAAR Surgical) in front of the natural lens for high myopia correction, without removing the natural lens.

\medterm{Intrathecal Pump (Baclofen Pump)}-- An implantable pump that delivers baclofen directly into the intrathecal space for the treatment of severe spasticity. The Medtronic SynchroMed II is the only FDA-approved implantable pump for intrathecal baclofen (ITB). A catheter is tunneled subcutaneously from the pump (placed in the lower abdominal wall) to the lumbar spine and threaded into the intrathecal space (typically at L2-L3). The pump reservoir (20 mL or 40 mL) is refilled percutaneously every 1-3 months through the central fill port. Indications: severe spasticity refractory to oral therapy (max dose of oral baclofen: 80 mg/day) due to multiple sclerosis, spinal cord injury (most common indication), cerebral palsy, traumatic brain injury, and stroke. Efficacy: reduces spasticity (Ashworth score decrease of 1-2 points) and improves function/pain. Screening test: 50-100 µg intrathecal bolus to assess response before implantation. Baclofen concentration: 500-2000 µg/mL. Flow rate: adjustable (continuous infusion, simple or complex dosing). Contraindications: active infection, small body habitus, and coagulopathy. Complications: pump pocket infection (most common, 5-10\%), catheter kink/occlusion/dislodgement, pump malfunction, baclofen overdose (respiratory depression, hypotonia, sedation — treat with physostigmine), baclofen withdrawal (severe rebound spasticity, fever, confusion, rhabdomyolysis — treat with oral baclofen/benzodiazepines), and CSF leak. Catheter-related failures account for 50\% of ITB system revisions.

\medterm{Intrauterine Device (IUD)}-- A small T-shaped device inserted into the uterine cavity for long-acting reversible contraception (LARC). Hormonal IUD (levonorgestrel-releasing): Mirena (52 mg, approved 8 years), Kyleena (19.5 mg, 5 years), Skyla (13.5 mg, 3 years). Copper IUD (non-hormonal): Paragard (T380A, approved 10-12 years, contains copper wire on arms and stem). Mechanism: hormonal IUD thickens cervical mucus, suppresses endometrial proliferation, and may inhibit ovulation; copper IUD creates a sterile inflammatory reaction toxic to sperm and ova. Insertion: after sterile speculum exam, uterine sound, the IUD is placed via inserter tube; strings trimmed to 2-3 cm from the external cervical os. Best time: during menses (cervix slightly open, pregnancy excluded). Pain management: NSAIDs, para/paracervical block (for nulliparous, highly anxious). Efficacy: >99\% (typical use failure rate <1\%), same as tubal ligation and vasectomy. Contraindications: pregnancy, current PID, purulent cervicitis, Wilson disease (copper IUD), uterine anomaly, and unexplained vaginal bleeding. Complications: perforation (1:1000, usually at insertion), expulsion (2-10\%, most in the first year), cramping, irregular bleeding (heavy menses with copper, spotting with hormonal), string problems (lost or retracted), and PID (primarily in the first 3 weeks post-insertion; higher risk with STI exposure). Return to fertility: immediate upon removal. Non-contraceptive benefits: reduces heavy menstrual bleeding (hormonal), endometrial cancer protection (hormonal), and emergency contraception (copper, up to 5 days after unprotected sex).

\medterm{Mastopexy (Breast Lift Implant)}-- A surgical procedure to lift and reshape sagging breasts (ptosis). Often combined with breast augmentation (implants) or reduction. Indications: involutional ptosis after pregnancy/breastfeeding, weight loss, aging. The NAC (nipple-areolar complex) is repositioned superiorly. The augmentation-mastopexy is a complex operation with a higher revision rate than augmentation alone.

\medterm{Mechanical Heart Valve}-- A prosthetic heart valve made of pyrolytic carbon (strong, thromboresistant, durable). The St. Jude Medical (now Abbott) bileaflet valve is the most common: two semicircular hinged leaflets that open centrally, creating three flow orifices with minimal obstruction. The Medtronic-Hall tilting disc valve (a single tilting disc suspended by a central strut) is less common. The Omnicarbon valve is another single-tilting disc. The Starr-Edwards ball-in-cage valve (caged-ball) is historical (no longer implanted but sometimes still seen in elderly patients). Advantage: excellent durability (>30 years). Disadvantage: lifelong therapeutic anticoagulation with warfarin (INR 2.0-3.0 for aortic, 2.5-3.5 for mitral) due to thrombogenicity. Risk of thromboembolism (1-3\%/year with adequate anticoagulation) and bleeding (1-3\%/year). Indications: age <50-60 (especially in the mitral position), kidney disease (bioprosthetic valves degenerate faster), and patient preference (willing to take warfarin). Sound: a distinct clicking/ticking sound audible to the patient and sometimes to others.

\medterm{Mid-Urethral Sling (MUS)}-- A polypropylene mesh tape placed under the mid-urethra to treat stress urinary incontinence (SUI) by providing dynamic urethral support during increases in intra-abdominal pressure. Retropubic (TVT, tension-free vaginal tape): a monofilament polypropylene mesh tape inserted via a vaginal incision and passed through the retropubic space with trocars exiting the suprapubic skin (risk of bladder perforation, bowel injury). Transobturator (TOT, TVT-O): the tape is passed through the obturator foramen exiting the groin crease (less bladder injury, more groin pain). Single-incision mini-sling (SIMS, e.g., Altis, MiniArc): shorter, the tape is anchored directly to the obturator internus fascia through a single vaginal incision (less invasive, but higher failure rates for severe SUI). Indications: urodynamic stress incontinence (USI, SUI confirmed on urodynamic studies) with hypermobility of the urethra. Contraindications: pregnancy (future childbearing), active UTI, untreated UTI, and desire to void. Complications: mesh erosion into the vagina (3-5\%), voiding dysfunction (retention, de novo urgency), groin/thigh pain (TOT), bladder perforation (TVT, 2-5\%), mesh infection, and de novo urgency urinary incontinence. Success rate: 80-90\% subjective cure at 5 years. The FDA 2011 transvaginal mesh safety communication specifically warned against mesh for POP repair; slings remain the standard of care for SUI.

\medterm{Pacemaker (Cardiac Implantable Electronic Device CIED)}-- An implantable device that delivers electrical pulses to the heart to maintain adequate heart rate. Single-chamber: VVI (RV lead, for rate support in AF). Dual-chamber: DDD (RA + RV leads, for AV block and sinus node dysfunction). Biventricular (CRT-P): RV + LV (coronary sinus lead) for cardiac resynchronization in heart failure. Indications: symptomatic sinus node dysfunction (sick sinus syndrome, HR <40), third-degree (complete) AV block, Mobitz II second-degree AV block, bifascicular block with intermittent AV block, and chronotropic incompetence. Pacing modes (NBG code): first letter = chamber paced (A, V, D), second = chamber sensed, third = response (I, T, D), fourth = rate modulation (R). Rate-responsive (DDDR): uses an accelerometer or minute ventilation sensor to increase rate during exercise. MRI-conditional (most modern): labeled MR Conditional, safe under specified conditions. Endocardial leads: passive fixation (tines) or active fixation (screw-in helix), placed via subclavian or axillary venous access. Left subpectoral pocket is standard. Leadless pacemaker (Micra AV, Medtronic; Aveir, Abbott): all-in-one device implanted in the RV via a femoral approach (no pocket, no lead). Dual-chamber leadless pacing is emerging (Aveir DR). Generator replacement: every 6-12 years (CRT-P shorter). Remote monitoring: standard (MyCareLink, Merlin, Latitude) reduces clinic visits. Complications: infection (pocket, endocarditis, 1-2\% for new implant, 5-7\% for system revision), lead failure (fracture, insulation break, recall-prone), pneumothorax (subclavian puncture), pocket hematoma, and twiddler syndrome (patient rotates generator, dislodging leads).

\medterm{Patent Ductus Arteriosus (PDA) Occlusion Device}-- An implant for transcatheter closure of a patent ductus arteriosus, a congenital heart defect where the fetal ductus arteriosus fails to close after birth. The Amplatzer Duct Occluder II (Abbott) and the Cook detachable coil are the most common. The device is delivered through a 4-6 Fr sheath via the femoral vein or artery. The Amplatzer device is a self-expanding nitinol mesh with a retention disc on the aortic side and a plug on the pulmonary side. Indications: hemodynamically significant PDA (left-to-right shunt, enlarged left heart, continuous murmur). Closure is typically performed >6 months of age (weight >6 kg). Absolute contraindication: pulmonary vascular disease (Eisenmenger syndrome with right-to-left shunt, PDA-dependent pulmonary/systemic circulation). Complications: device embolization (most common), residual shunt, hemolysis, and obstruction of the LPA or descending aorta.

\medterm{Penile Prosthesis (Implant)}-- A surgically implanted device for treating erectile dysfunction (ED) refractory to oral therapy (PDE5 inhibitors) and vacuum/ICI therapies. Malleable (semi-rigid) rods: AMS 600M/650M (Boston Scientific), Coloplast Genesis. These are always semirigid and can be bent downward for urination/intercourse; simpler and cheaper, but less natural. Inflatable prostheses (two-piece: AMS Ambicor; three-piece: AMS 700 (Boston Scientific) with LGX (length-expanding) or CX (controlled expansion) cylinders, Coloplast Titan with One-Touch Release): cylinders placed in the corpora cavernosa, a scrotal pump (in the scrotum), and a fluid reservoir (retropubic space, collapsible for Coloplast; pre-connected for AMS 700). Three-piece is the standard for the best results (natural erection and flaccidity). Indications: organic ED (diabetes, post-radical prostatectomy, Peyronie disease, severe vascular disease) refractory to medical therapy. Contraindications: active infection, untreated UTI, and severe penile fibrosis (scarring). Surgery: penoscrotal (most common) or infrapubic approach. Cylinder size is measured intraoperatively (dilatation with Hegar dilators, Furlow inserter). Complication rates: infection (1-3\% in naive patients; higher in revision/revision, 5-15\%), mechanical failure (5-10\% at 5 years, higher for 3-piece), cylinder erosion, pump migration, and auto-inflation (Correction: the Coloplast Titan has a lock-out valve to prevent auto-inflation). Satisfaction rates: >90\% (patient and partner). Antibiotic impregnation: InhibiZone (rifampin/minocycline) on AMS devices reduces infection.

\medterm{Port-a-Cath (Implantable Venous Access Port)}-- A fully implanted central venous access device consisting of a silicone catheter (inserted into the superior vena cava via the internal jugular or subclavian vein) connected to a titanium or plastic reservoir (port) placed in a subcutaneous pocket on the chest wall (typically 2-3 cm below the clavicle, over the pectoralis major fascia). Access is through the skin via a non-coring Huber needle (right-angle or straight). The power-injectable port (e.g., Bard PowerPort, Cook Vital-Port) can handle CT contrast injection at up to 5 mL/sec. Indications: chemotherapy (most common), long-term IV antibiotics (osteomyelitis, endocarditis), TPN, frequent blood transfusions, and frequent blood draws (rare). Double-lumen ports are available for incompatible infusions. Placement: interventional radiology or surgeon, under ultrasound and fluoroscopic guidance. Tip position: cavoatrial junction. Flush: heparinized saline (100 U/mL, 5 mL) monthly when not in use, or saline-only ports (e.g., BD Q-Syte, neutral displacement). Complication rate: 5-15\%. Immediate: pneumothorax, hemothorax, arterial puncture, air embolism. Short-term: pin-hole infection, catheter malposition. Long-term: catheter-related bloodstream infection (CRBSI, most common, 1-3 per 1000 catheter-days, treat with antibiotics + port removal if tunnel infection), catheter occlusion (thrombotic: alteplase), fibrin sheath formation (retraction of blood, resistant aspiration but easy flushing), venous thrombosis (upper extremity DVT, catheter-related thrombosis), and port erosion/exposure. Port longevity: indefinite (until no longer needed or complication occurs).

\medterm{Responsive Neurostimulation (RNS System)}-- An implantable closed-loop neurostimulation system for treating drug-resistant focal epilepsy. The NeuroPace RNS System is the only FDA-approved device. Components: a cranially implanted neurostimulator (placed in a skull defect, similar to a cranial bone flap), one or two depth or subdural strip electrodes placed at the seizure focus(es), and a wireless telemetry wand for programming and data download. The device continuously records electrocorticography (ECoG) data and delivers stimulation (brief electrical pulses) when it detects a pattern that matches a pre-programmed seizure-onset pattern (detection and stimulation cycles). It is a responsive, adaptive, closed-loop system (not continuous like DBS). Indications: adults with drug-resistant focal epilepsy (with 1-2 seizure foci) who are not candidates for resective surgery. Contraindications: multifocal epilepsy with >2 foci, generalized epilepsy, and MRI incompatibility (RNS is MRI-conditional). Efficacy: 50-70\% median seizure reduction at 2 years, continues to improve over time (possibly due to neuromodulation effects). Advantages over VNS and DBS: delivers therapy only when needed, provides long-term ECoG monitoring for seizure tracking, and is adjustable without surgery. Complications: infection (5-10\%, often requires device removal), implant site pain, hemorrhage, and lead fracture. Device battery: about 4-6 years.

\medterm{Shunt (Ventriculoperitoneal VP Shunt)}-- A device that diverts cerebrospinal fluid (CSF) from the cerebral ventricles to the peritoneal cavity for the treatment of hydrocephalus. Components: ventricular catheter (inserted through a frontal or occipital burr hole into the lateral ventricle, typically the right frontal horn), a valve (programmable: Codman Hakim, Medtronic Strata, Sophysa Polaris; or fixed-pressure: low/medium/high), and a peritoneal catheter (tunneled subcutaneously from the head to the upper abdomen, inserted into the peritoneal cavity). Programmable valves: adjusted externally with a magnetic tool (Codman Hakim: 30-200 mm H$_2$O; Medtronic Strata: 0.5-2.5; Aesculap Miethke proGAV 2.0: gravitational valve that changes with body position). Indications: congenital hydrocephalus, normal pressure hydrocephalus (NPH), post-hemorrhagic hydrocephalus (IVH, SAH), post-meningitic hydrocephalus, and obstructive hydrocephalus from tumors. Contraindications: active peritonitis, infected CSF (VP shunt), and presence of a ventriculostomy (EVD) that is infected. Complications: shunt obstruction (most common, 30-40\% at 2 years, most frequent in the proximal catheter), infection (5-10\%, most within 3 months, typically Staphylococcus epidermidis), overdrainage (subdural hematoma, slit ventricle syndrome, orthostatic headaches), underdrainage (persistent hydrocephalus), and shunt migration (peritoneal catheter can erode into the bowel, bladder, or through the skin). Overdrainage in NPH may lead to subdural hygromas/hematomas from negative pressure. Antibiotic-impregnated catheters (Bactiseal, Codman) reduce infection (OR 0.5-0.6).

\medterm{Spinal Cord Stimulator (SCS)}-- An implantable neuromodulation device for treating chronic neuropathic pain. Components: epidural leads (percutaneous cylindrical leads or surgical paddle leads placed in the dorsal epidural space) and an implantable pulse generator (IPG, battery) placed in the lower back/buttock/subcostal pocket. Trial: temporary external system for 3-7 days to evaluate efficacy before permanent implantation (>50\% pain reduction needed for proceeding). Indications: failed back surgery syndrome (FBSS, most common), complex regional pain syndrome (CRPS, type I and II), diabetic peripheral neuropathy, and refractory angina. Traditional SCS: paresthesia-based (tonic, 40-60 Hz) covers the painful area with a tingling sensation (paresthesia overlap). Newer waveforms: high-frequency (10 kHz, HF10 therapy) — paresthesia-free, superior for back pain; burst stimulation (40 Hz burst with 5 spikes) — targets affective component of pain; and closed-loop SCS (Abbott Proclaim, closed-loop feedback from evoked compound action potentials ECAPs) — maintains consistent activation and improves outcomes. Brands: Boston Scientific (Spectra WaveWriter, paresthesia-based + burst + combination), Abbott (Proclaim DRG, for SCS and dorsal root ganglion stimulation), Nevro (Senza HFX, high-frequency 10 kHz), and Medtronic (Intellis with differential target multiplexed programming). DRG stimulation (Abbott Axium): lead placed at the dorsal root ganglion in the epidural space at the L1-S2 levels for focal neuropathic pain (e.g., post-herniorrhaphy pain, CRPS of a limb) with better paresthesia coverage and positional stability. Complications: lead migration (most common, 5-15\%), infection (2-5\%, often IPG pocket), IPG site pain, lead fracture, dural puncture, and CSF leak.

\medterm{Subdermal Contraceptive Implant (Nexplanon)}-- A single-rod, non-biodegradable, etonogestrel-releasing implant (68 mg, released over 3 years) placed subdermally on the inner upper arm. The applicator is pre-loaded (loaded disposable inserter, changed from Implanon). Mechanism: inhibits ovulation (by suppressing LH surge), thickens cervical mucus, and suppresses endometrial proliferation. Insertion: after local anesthesia—the device is inserted subdermally at the medial aspect of the non-dominant arm, 8-10 cm above the medial epicondyle, in the sulcus between the biceps and triceps. The new Nexplanon uses a single-use, preloaded, auto-disable applicator (prevents reinsertion). The device is radiopaque (barium sulfate added). Efficacy: >99\% (Pearl Index 0.05). Duration: 3 years (FDA-approved), but evidence suggests it is effective up to 5 years. Advantages: long-acting reversible contraception (LARC), 100\% reversible (returns to fertility immediately after removal), excellent for adolescents and women who cannot use estrogen. Disadvantages: irregular bleeding/spotting (most common side effect, 50\% after 1 year), amenorrhea (20\%), weight gain, breast tenderness, and acne. Contraindications: active VTE, liver disease, breast cancer, and unexplained abnormal vaginal bleeding.

\medterm{Transcatheter Aortic Valve Replacement (TAVR Valve)}-- A bioprosthetic valve crimped onto a catheter and delivered percutaneously to replace a stenotic aortic valve without open heart surgery. Two main platforms: balloon-expandable (Edwards SAPIEN 3 Ultra, SAPIEN 3 — bovine pericardial leaflets, cobalt-chromium stent frame) and self-expanding (Medtronic Evolut PRO+, PRO — porcine pericardial leaflets, nitinol frame with external pericardial wrap for sealing). Access: transfemoral (most common), transapical (left lateral thoracotomy, direct LV apex), transaortic (direct ascending aorta), transaxillary/subclavian, and transcarotid. Indications: symptomatic severe aortic stenosis (reduced survival if untreated) — currently approved for all risk categories (extreme, high, intermediate, and low risk). The PARTNER trials (Edwards) and CoreValve US trials (Medtronic) validated TAVR across risk strata. Pre-procedural CT: annulus sizing (perimeter, area), coronary height (risk of coronary occlusion), and iliofemoral access assessment. Annulus measurement is crucial: oversizing by 5-15\% prevents paravalvular leak; severe oversizing risks annular rupture. Valve-in-valve: TAVR within a failed surgical bioprosthesis (viable TAVR) is safe and effective. Complications: paravalvular leak (PVL, 5-10\% moderate or greater for new-generation valves, less with improved sealing skirts), permanent pacemaker requirement (PPM, 15-30\% for self-expanding — due to conduction disturbances LBBB/AV block, 5-10\% for balloon-expandable), stroke (2-5\%, lowered with embolic protection devices like Sentinel CPS), vascular access complications (bleeding, pseudoaneurysm, dissection, 5-10\%), coronary obstruction (rare but fatal if undetected), and annular rupture (rare).

\medterm{Transjugular Intrahepatic Portosystemic Shunt (TIPS)--See TIPS in Surgery.}-- See also Intrathecal Pump for CSF shunting.

\medterm{Tympanostomy Tube (Grommet)}-- A small ventilation tube inserted through the tympanic membrane (eardrum) to equalize pressure and drain middle ear fluid. Materials: silicone, fluoroplastic, and titanium; with or without antibiotic coating. Types: short-term (Paparella type 1, Armstrong, Shepard, Donaldson — extrude spontaneously in 6-18 months) and long-term (T-tube, Goode T-tube — stay in 2-4+ years, may require surgical removal). Indications: persistent otitis media with effusion (OME, >3 months in children with hearing loss >20 dB), recurrent acute otitis media (AOM, >3 episodes/6 months or >4/12 months), chronic eustachian tube dysfunction, and adult barotrauma (flyers, divers). Contraindications: active otorrhea, cholesteatoma, and prior history of TM perforation. Procedure: myringotomy (anterior-inferior quadrant, avoiding the chorda tympani nerve) with aspiration of effusion, then tube insertion. Follow-up: keep water out of ears (earplugs for swimming are controversial; generally recommended to avoid deep water diving but surface swimming is acceptable). Tubes extrude spontaneously as the TM regenerates. Complications: tube otorrhea (most common, 10-30\%, treat with topical antibiotics), obstruction (crust, dried effusion), early extrusion, persistent TM perforation after tube extrusion (1-5\%, higher with T-tubes), tympanosclerosis (TM scarring, 20-30\%), and cholesteatoma.

\medterm{Vagus Nerve Stimulator (VNS)}-- An implantable neurostimulator that delivers intermittent electrical pulses to the left vagus nerve for the treatment of drug-resistant epilepsy and depression. The LivaNova VNS Therapy System (Demipulse, AspireHC) is the primary device. A spiral electrode is wrapped around the left vagus nerve in the carotid sheath (left is chosen to minimize cardiac effects, as the right vagus heavily innervates the SA node). The pulse generator is placed in a left infraclavicular subcutaneous pocket. Stimulation parameters: output current (0.25-3.0 mA), frequency (20-30 Hz), pulse width (250-500 µs), on-time (30 sec), and off-time (5 minutes). Indications: drug-resistant partial-onset epilepsy (\(\ge\)12 years, FDA-approved 1997) and treatment-resistant depression (FDA-approved 2005). It is not for patients with prior left cervical vagotomy. Efficacy: 30-50\% seizure reduction at 3 months, improving to 50-60\% at 2 years. Magnets: patients/caregivers can swipe a magnet over the generator to trigger on-demand stimulation to abort seizures. Vagal nerve stimulation for depression: NIMH STAR*D post-hoc analysis showing benefit, but insurance approval varies. Complications: voice alteration/hoarseness (most common, 30-60\%, due to stimulation of the laryngeal nerve, resolves when stimulator off), cough, dysphagia, paresthesias, infection, and device malfunction. It is MRI conditional. Battery life: 4-10 years.

\medterm{Vascular Graft (Synthetic)}-- A prosthetic conduit used to replace or bypass diseased arteries or veins. Materials: Dacron (polyethylene terephthalate PET, woven or knitted, used for large-diameter high-flow vessels: aorta, iliac) and ePTFE (expanded polytetrafluoroethylene, Gore-Tex, used for medium-diameter vessels: femoral, popliteal, AV grafts for dialysis). Heparin-bonded grafts (e.g., Gore Propaten, GORE-TEX with heparin bioactive surface) reduce early thrombosis. Linings: tissue-engineered and autologous endothelial cell-seeded grafts (experimental). Bifurcated grafts: for aortoiliac occlusion or AAA repair (body + left/right limbs). Ringed grafts: external ePTFE rings prevent kinking at the knee (fem-pop bypass). Biological grafts: autologous (saphenous vein — best patency for infrainguinal bypass, 60-80\% at 5 years ; internal mammary artery for CABG), allograft (cadaveric vein), and xenograft (bovine carotid artery, bovine mesenteric vein). Indications: peripheral arterial disease (aortobifemoral, femoropopliteal, femorotibial bypass), aortic aneurysm repair (open AAA repair or hybrid repair), and hemodialysis access (AV graft when native AV fistula is not feasible, e.g., basilic vein transposition, forearm loop). Complications: infection (most dreaded, especially if graft is in the groin, treatment includes excision and extra-anatomic bypass), thrombosis (pseudointimal hyperplasia, especially at low-flow anastomoses), anastomotic pseudoaneurysm, and graft-enteric fistula (aortoduodenal fistula after AAA repair). Patency: aortobifemoral 80-90\% at 5 years; femoropopliteal 50-70\%.

\medterm{Vena Cava Filter (IVC Filter)}-- An endovascular device placed in the inferior vena cava to prevent pulmonary embolism (PE) in patients with venous thromboembolism (VTE) who cannot tolerate anticoagulation. The filter is a metal cage (cone-shaped, with struts and hooks) that traps emboli while allowing blood flow. The most commonly used is the Greenfield filter (stainless steel or titanium, 4-6 struts). Retrievable filters: Bard Denali (conical, retrievable, nitinol), Cook Celect (conical, retrievable, with timers), Cook Gunther Tulip (conical, retrievable), and Cordis OptEase (double-basket, retrievable). Permanent filters: Greenfield (stainless steel) and TrapEase. Indications: contraindication to anticoagulation (active bleeding, recent surgery, hemorrhagic stroke), failure of anticoagulation (recurrent PE despite therapeutic anticoagulation), and prophylactic in high-risk trauma (severe pelvic/long bone fractures, TBI). Contraindications: complete IVC thrombosis (no room to deploy), severe IVC stenosis, and uncorrectable coagulopathy. The ideal position is infrarenal IVC (below the renal veins, above the IVC bifurcation at L2-L3). The filter is placed under fluoroscopic guidance via the right femoral vein (most common) or right internal jugular vein. Complications: IVC penetration (struts through the IVC wall, estimated 15-40\% in retrievable filters, often asymptomatic), filter migration (to the heart, pulmonary artery — can be fatal), IVC thrombosis (5-15\%, causing phlegmasia, edema, DVT), retrieval failure (filter tilting, endothelialization, strut penetration, significant clot in filter, time since implant >1 year reduces likelihood of successful retrieval — 60-90\% retrieval success in experienced centers). DVT at the insertion site (5-10\%). The FDA in 2010 and 2014 recommended that retrievable filters be removed as soon as PE protection is no longer needed (ideally within 29-54 days). Long-term complications: filter fracture, systemic embolization, and chronic IVC occlusion.

\medterm{Ventricular Assist Device (VAD/LVAD)}-- A mechanical pump that supports the failing left ventricle. The HeartMate 3 (Abbott) is the current state-of-the-art LVAD: a fully magnetically levitated centrifugal pump with artificial pulse (designed to reduce pump thrombosis and shear stress). The HeartWare HVAD (Medtronic, previously HeartWare — withdrawn from market due to mortality/morbidity concern but still used in some patients) is a smaller centrifugal pump with a hybrid suspension. Impella (Abiomed): a microaxial percutaneous VAD placed via the femoral artery (Impella CP, Impella 5.5) for short-term support (up to 30 days for Impella 5.5). The TandemHeart (LivaNova) is a percutaneous left atrial-to-femoral artery bypass system. Indications: bridge to transplant (BTT, keep patient alive until a donor heart becomes available), bridge to candidacy (BTC, improve end-organ function to become transplant-eligible), destination therapy (DT, permanent support for patients ineligible for heart transplant), and bridge to recovery (BTR, support until the native heart recovers, e.g., myocarditis, post-cardiotomy). Surgery: median sternotomy or thoracotomy, inflow cannula in the LV apex, outflow graft to the ascending aorta (LVAD). A driveline exits the abdominal wall and connects to an external controller and batteries. Complications: pump thrombosis (reduced with HeartMate 3 full mag-lev design), stroke (ischemic and hemorrhagic, 10-20\% at 1-2 years), driveline infection (15-30\% at 1 year, 30-50\% at 2 years), GI bleeding (30-50\% from acquired von Willebrand syndrome type 2A, especially with axial flow HM II), right heart failure (20-30\%), aortic insufficiency (from continuous flow, 20-30\% at 2-3 years), and hemolysis. Life: HM3 is approved for up to 2 years as DT but is known to run much longer.

